% !TEX TS-program = pdflatex
\documentclass[11pt,a4paper]{article}

% ---------- Minimal, GitHub-friendly setup ----------
\usepackage[utf8]{inputenc}   % for pdflatex
\usepackage[T1]{fontenc}
\usepackage{lmodern}
\usepackage[margin=1in]{geometry}
\usepackage{parskip}          % space between paragraphs, no indents
\usepackage{hyperref}         % clickable links
\usepackage{enumitem}         % compact lists
\usepackage{xcolor}
\usepackage{graphicx}

% ---------- Code listings (lightweight) ----------
\usepackage{listings}
\lstdefinestyle{ghpython}{
	language=Python,
	basicstyle=\ttfamily\small,
	numbers=left,
	numberstyle=\tiny,
	stepnumber=1,
	numbersep=8pt,
	showstringspaces=false,
	breaklines=true,
	frame=single,
	framerule=0.4pt,
	tabsize=2,
	upquote=true,
	keywordstyle=\color{blue!60!black}\bfseries,
	commentstyle=\color{gray!70},
	stringstyle=\color{green!40!black}
}
\lstset{style=ghpython}

% ---------- Safe inline code macro (handles underscores) ----------
\newcommand{\code}[1]{\texttt{\detokenize{#1}}}

% ---------- Handy helpers ----------
\newcommand{\func}[1]{\texttt{#1()}}
\newcommand{\class}[1]{\texttt{#1}}

% ---------- Metadata ----------
\title{\textbf{GA Sudoku Solver (PyGAD)}\\
	\large Minimal Documentation}
\author{Project Documentation}
\date{\today}

\begin{document}
	\maketitle

	\begin{center}
		\vspace{-0.5em}
		\small
		Repository: \href{github.com/ArmanBjr/Artificial_Intelligence-Course}{https://github.com/ArmanBjr/Artificial\_Intelligence-Course} \\
		Document version: \today
	\end{center}

	\hrule
	\vspace{0.8em}

	\begin{abstract}
		This document provides a concise, contributor-facing overview of the Sudoku solver implemented with a Genetic Algorithm (PyGAD). It focuses on the important public functions and how they fit together. The style and package set are kept lightweight for GitHub repositories.
	\end{abstract}

	\vspace{0.5em}
	\tableofcontents
	\vspace{0.8em}
	\hrule
	\vspace{1.2em}

	% =========================
	\section{Overall Overview}
	% =========================
	This project solves Sudoku puzzles using a Genetic Algorithm (GA) built on \code{pygad}. Only the \emph{empty} cells are represented as genes; original clues remain fixed. The main ideas:

	\begin{itemize}[leftmargin=1.2em]
		\item \textbf{Per-box initialization:} Each chromosome starts with every 3$\times$3 box filled by its missing digits (shuffled), giving box-valid candidates from the beginning.
		\item \textbf{Box-preserving operators:}
		\begin{itemize}
			\item \emph{Crossover} (\func{custom\_crossover}): offspring inherit entire 3$\times$3 boxes from parents.
			\item \emph{Mutation} (\func{custom\_mutation}): swaps two genes within the same box, preserving the box’s multiset of digits.
		\end{itemize}
		\item \textbf{Fitness} (\func{fitness\_function}): penalizes duplicates across rows, columns, and boxes. The GA maximizes the negative penalty; a valid Sudoku yields fitness $=0$.
		\item \textbf{Solver} (\func{solve}): configures and runs the GA (population, generations, selection, crossover/mutation hooks), then decodes the best chromosome back into a grid.
	\end{itemize}

	This box-aware design constrains search to mostly fixing row/column conflicts, improving stability and time-to-solution.

	% =========================
	\section{Quick Start }
	% =========================
	\begin{enumerate}[leftmargin=1.2em]
		\item Install dependencies: \code{pygad}, \code{numpy}.
		\item Choose a puzzle from \code{SUDOKU_EXAMPLES}.
		\item Run \code{__main__} and inspect the output grid and fitness.
	\end{enumerate}

	% =========================
	\section{Repository Structure }
	% =========================
	\begin{verbatim}
			.
			|-- session4/
			|   |-- Sudoku.py
			|   |-- sample/
			|   |-- sudoku1.pdf
			|   `-- sudoku2.pdf
		\end{verbatim}


	% =========================
	\section{Key Components }
	% =========================
	\subsection{Class \class{SudokuSolver}}
	\begin{itemize}[leftmargin=1.2em]
		\item \func{solve} \dotfill GA configuration and execution (returns best grid and fitness).
		\item \func{fitness\_function} \dotfill Negative penalty based on duplicates in rows, columns, and boxes.
		\item \func{custom\_crossover} \dotfill Box-wise inheritance from parents.
		\item \func{custom\_mutation} \dotfill In-box gene swap operator.
	\end{itemize}

	\subsection{Utilities}
	\begin{itemize}[leftmargin=1.2em]
		\item \code{decode_solution} \dotfill Map chromosome back to a 9$\times$9 grid.
		\item \code{print_sudoku}, \code{analyze_solution} \dotfill I/O and simple validity checks.
	\end{itemize}

	% =========================
	\section{Example Snippet }
	% =========================
	\begin{lstlisting}
		solver = SudokuSolver(SUDOKU_EXAMPLES["expert"])
		solution, fitness = solver.solve(num_generations=1000,
		num_parents_mating=80,
		sol_per_pop=400)
	\end{lstlisting}

\end{document}
